\begin{abstractpage}

\begin{abstract}{romanian}
Lucrarea de licență prezintă proiectarea și implementarea unei aplicații web pentru planificarea itinerariilor turistice, realizată cu ASP.NET Core MVC și Entity Framework Core. Sistemul permite utilizatorilor autentificați să creeze călătorii, să adauge locații și rezervări, să gestioneze tipuri de obiective prin etichete și să genereze automat planuri zilnice de vizitare.

Generatorul de itinerariu construiește agenda fiecărei zile pe baza datei de început și de sfârșit, orei de explorare, rezervărilor existente și duratei medii de vizitare pentru fiecare locație. Algoritmul prioritizează întâi evenimentele programate, apoi adaugă obiective disponibile din listă în ordinea proximității și a compatibilității cu intervalele orare. Lucrarea detaliază construcția modelelor de date relaționale, configurarea entităților many-to-many, validarea datelor și optimizarea fluxului de creare a itinerariului.

Rezultatul este o aplicație care sprijină utilizatorii în organizarea eficientă a traseelor turistice, sinteza rezervărilor și planificarea responsabilă a zilelor de vizitare, oferind o experiență coerentă pe tot parcursul procesului.
\end{abstract}

% \begin{abstract}{english}
% This bachelor's thesis presents the design and implementation of a web application for travel itinerary planning, built with ASP.NET Core MVC and Entity Framework Core. The system enables authenticated users to create trips, add locations and reservations, manage categories of attractions through tags, and automatically generate daily visit plans.

% The itinerary generator builds each day's schedule based on trip dates, exploration hours, existing reservations, and the average visit duration for each location. The algorithm prioritizes scheduled events first, then adds available attractions by proximity and time window compatibility. The thesis details the construction of relational data models, many-to-many entity configuration, input validation, and itinerary generation flow optimization.

% The result is an application that supports users in organizing travel routes efficiently, combining bookings and daily planning into a cohesive experience throughout the trip preparation process.
% \end{abstract}

\end{abstractpage}